\chapter{Anexo IV: Manual de Usuario}
\label{cap:Manual}
En este anexo se incluye un manual de usuario para la solución desarrollada. Este manual está dirigido a los usuarios finales que interactuarán con el sistema, proporcionando instrucciones claras y detalladas sobre cómo utilizar las funcionalidades implementadas. Se incluyen capturas de pantalla, descripciones de cada función y pasos para realizar tareas comunes dentro del sistema.

La base física de esta jerarquía se fundamenta en el siguiente elemento:\par

\subsection{Nivel 0: Marco Base (\texttt{00\_Main\_Interface})}
Representa la interfaz matriz o Nivel 0 del SCADA (Figura \ref{fig:00_main_interface}). En lugar de ejecutar transiciones de pantalla completa, esta plantilla actúa como un contenedor estático  que alberga el menú de navegación principal e implementa dos \textit{Picture Windows} independientes.\par

La ventana de trabajo central ocupa la mayor parte de la pantalla y constituye el área de visualización dinámica donde se cargan todas las interfaces operativas (Niveles 1 a 4). Paralelamente, la ventana de alarmas inferiores es un espacio de menores dimensiones anclado de forma fija en el pie de página, cuya única función es garantizar la supervisión ininterrumpida de los eventos críticos independientemente de la profundidad de navegación del usuario.
\FloatBarrier
\subsection{Nivel 1: Supervisión Global}
\label{sec:nivel_1_overview}
\subsubsection{Vista General de aerogeneradores (\texttt{10\_WTG\_Overview})}
La pantalla \texttt{10\_WTG\_Overview} actúa como el nodo principal de enrutamiento dinámico del SCADA. Su diseño optimiza la navegación hacia los niveles inferiores mediante tres mecanismos clave integrados en su cabecera.\par

Por un lado, se ha implementado un sistema de filtrado mediante dos selectores desplegables (Figura \ref{fig:interfaz_grafica_selector}) que acotan la información visualizada, actualizando los aerogeneradores mostrados sin necesidad de realizar transiciones a otras pantallas ( Algoritmo \ref{alg:filtrado_dinamico}, vease la función \texttt{AplicarFiltroPosicion()} en el Anexo \ref{cap:anexo_codigo_solucion}).\par

Por otro lado, como se observa en la figura \ref{fig:interfaz_grafica_park} un selector central gestiona el archivo alojado en el \textit{Picture Window} de trabajo. Permite alternar instantáneamente entre la telemetría individual de las turbinas (\texttt{20\_WTG\_Turbine\_Overview.pdl}) y la vista de datos por parques (\texttt{21\_WTG\_Park\_Overview.pdl}).\par

Finalmente, como se muestra en la figura \ref{fig:interfaz_grafica} una matriz superior de indicadores discretos organizados por parques representa el estado de cada aerogenerador. Seleccionar uno de ellos actúa como acceso directo al Nivel 4 (\texttt{40\_WTG\_Gen}), omitiendo pasos intermedios e inyectando automáticamente el \textit{Tag Prefix} correspondiente para cargar los datos de la máquina seleccionada.\par


\label{sec:nivel_1_subestacion}
Esta pantalla actúa como un contenedor especializado para la visualización de la infraestructura eléctrica del parque. Arquitectónicamente, se ha diseñado como una interfaz intermedia que alberga un único \textit{Picture Window} destinado a cargar los diversos diagramas unifilares (\textit{Single Line Diagrams}).\par

La decisión de no integrar estos esquemas directamente en la pantalla raíz (\texttt{00\_Main\_Interface}) y emplear este encapsulamiento responde a dos criterios técnicos fundamentales, por un lado, la gestión de la visualización y desplazamiento de diagramas eléctricos complejos (\textit{scrolling}), y por otro, el aislamiento de variables mediante \textit{Tag Prefixes} para evitar interferencias con el entorno principal del SCADA.\par
\subsubsection{Interfaz de monitorización, alarmas y tendencias (\texttt{12\_MET\_Station}, \texttt{13\_SYS\_Alarms} y \texttt{14\_SYS\_Trends})}

Como complemento indispensable a las vistas generales de generación y evacuación de potencia, se han integrado en el Nivel 1 tres interfaces dedicadas a la monitorización ambiental, el diagnóstico del historial de fallos y el soporte operativo del entorno SCADA. Estas pantallas, identificadas bajo los archivos \texttt{12\_MET\_Station.pdl}, \texttt{13\_SYS\_Alarms.pdl} y \texttt{14\_SYS\_Trends\_Interface.pdl}, aseguran que el operador mantenga una conciencia situacional periférica sobre las variables externas, las condiciones meteorológicas y el estado de salud del propio sistema de adquisición de datos sin necesidad de descender individualmente en la jerarquía de cada aerogenerador.\par

\FloatBarrier

\section{Interfaz Gráfica: Nivel 2 - Supervisión Global por Activos}
\label{sec:nivel2_scada}

Habiendo definido en el Nivel 1 las áreas macroscópicas y los marcos contenedores del sistema \cite{7.6.3}, el Nivel 2 se constituye como el primer escalón de desglose y segregación funcional del parque eólico. Siguiendo de forma estricta el patrón de diseño e ingeniería de nomenclatura definido en la Ecuación~\ref{eq:nomenclatura}, todas las pantallas pertenecientes a este nivel adoptan el prefijo numérico \texttt{2X}.

\subsection{Definición y Catálogo de Pantallas del Nivel 2}
Con el propósito de gestionar de forma eficiente la telemetría de los 120 aerogeneradores y el flujo energético de las 9 subestaciones de transformación y evacuación \cite{7.3}, se han diseñado e implementado las siguientes interfaces gráficas en formato \texttt{.pdl} dentro del entorno de desarrollo SIMATIC WinCC V8.1 \cite{29}:

\begin{itemize}
    \item \textbf{\texttt{20\_WTG\_Turbine\_Overview.pdl} (Malla de Turbinas Individuales):} Pantalla orientada a objetos que actúa como un contenedor dinámico en tiempo de ejecución (\textit{Runtime}) \cite{575}. Alberga un total de 120 ventanas de imagen (\textit{Picture Windows}) indexadas numéricamente (\texttt{PW\_1} a \texttt{PW\_120}) \cite{1748, 1763}. Su propósito operativo es renderizar de forma compacta y ordenada los \textit{faceplates} de monitorización rápida de cada activo físico en base al estado de los filtros globales de la aplicación.
    \item \textbf{\texttt{21\_WTG\_Park\_Overview.pdl} (Supervisión Agregada por Clúster):} Interfaz gráfica de abstracción técnica superior orientada a la consolidación de datos e indicadores clave de rendimiento (KPIs) por subestaciones o zonas lógicas \cite{1167, 1279}. Esta vista unifica variables totalizadas del clúster (potencia activa agregada, velocidad media del viento sectorial y balance de activos en estado de alarma o mantenimiento) \cite{1241, 1280}.
    \item \textbf{\texttt{22\_SUB\_SLD\_General.pdl} (Diagrama Unifilar General del Complejo):} Representación estática y dinámica del esquema unifilar (\textit{Single Line Diagram} - SLD) macroscópico del parque eólico \cite{1169, 1200}. Visualiza de forma unificada las líneas de barras comunes del complejo y el estado de los interruptores automáticos (\texttt{XCBR}) y seccionadores (\texttt{XSWI}) principales de evacuación hacia la red de transporte nacional eléctrica \cite{855, 887}.
    \item \textbf{\texttt{23\_SYS\_TrendTime.pdl} (Tendencias Temporales de Proceso):} Pantalla especializada que integra el control nativo \textit{WinCC Online Trend Control} \cite{1804}. Permite al operador analizar de forma gráfica series temporales continuas de variables de proceso críticas (e.g., velocidad angular del rotor \texttt{WROT.RotSpd} o temperatura del estator \texttt{WGEN.Tmp}) \cite{1064, 1103}.
    \item \textbf{\texttt{24\_SYS\_XYTrend.pdl} (Gráficas de Diagnóstico Analítico X-Y):} Interfaz dedicada a la monitorización predictiva del rendimiento aerodinámico \cite{362, 441}. Su función primordial es cruzar la velocidad instantánea del viento local medida por el anemómetro de góndola (\texttt{WMET.WndSpd}) frente a la potencia activa final inyectada en bornes (\texttt{WGEN.W}) \cite{788, 1848, 1856}. De este modo, se genera dinámicamente en pantalla la curva de potencia real del activo, permitiendo evaluar desviaciones frente a la curva teórica del fabricante.
\end{itemize}

\subsection{Lógica Matemática de Ordenación en Cuadrícula y Filtrado}
Para evitar que el operador sufra de saturación de información ante 120 activos simultáneos, la pantalla \texttt{20\_WTG\_Turbine\_Overview} ejecuta un algoritmo dinámico en lenguaje ANSI-C a través del subsistema \textit{Global Scripts} \cite{586, 1285}. El script lee los estados lógicos de los filtros seleccionados por el usuario y posiciona vectorialmente en una matriz bidimensional los faceplates válidos.

Sea $i \in \{1, 2, \dots, N_{total}\}$ el índice único que identifica a cada aerogenerador, donde $N_{total} = 120$ \cite{1751}. Definimos las variables de configuración espacial inicial en el lienzo gráfico como:
\begin{itemize}
    \item $X_{start} = 10 \text{ px}$ \quad (Margen izquierdo de la cuadrícula) \cite{1744}
    \item $Y_{start} = 0 \text{ px}$ \quad (Margen superior de la cuadrícula) \cite{1744}
    \item $\Delta X = 475 \text{ px}$ \quad (Ancho del faceplate más el espaciado horizontal) \cite{1745}
    \item $\Delta Y = 220 \text{ px}$ \quad (Alto del faceplate más el espaciado vertical) \cite{1745}
    \item $C_{max} = 4$ \quad (Número máximo de columnas permitidas por fila) \cite{1745}
\end{itemize}

Introduciendo la variable indexada de conteo de activos válidos que han superado con éxito las condiciones de filtrado cromo-operativo como $k$ (inicializada en $k = 0$), la posición espacial final de la esquina superior izquierda (coordenadas \textit{Left} y \textit{Top} de WinCC) para cada ventana de imagen visible se calcula de forma unívoca mediante las siguientes funciones de mapeo discreto:

\begin{equation}
    X_{pos}(k) = X_{start} + \left( k \pmod{C_{max}} \right) \cdot \Delta X
    \label{eq:pos_x}
\end{equation}

\begin{equation}
    Y_{pos}(k) = Y_{start} + \left\lfloor \frac{k}{C_{max}} \right\rfloor \cdot \Delta Y
    \label{eq:pos_y}
\end{equation}

Donde la operación módulo ($\pmod{}$) restringe el crecimiento en el eje horizontal asegurando la matriz de 4 columnas, y la función piso matemático ($\lfloor \rfloor$) incrementa de forma discreta la altura del renglón cada vez que el contador $k$ supera un múltiplo exacto de $C_{max}$ \cite{1769}. 

\subsection{Mecanismo de Integración en el Marco Base (Runtime)}
En concordancia con la arquitectura multipuesto de alta disponibilidad y virtualización industrial mediante el servicio SIVaaS implementado \cite{600, 648}, el flujo de navegación prohíbe las transiciones de pantalla completa tradicionales para mitigar la sobrecarga de intercambio de memoria en las estaciones cliente \cite{1179}. 

Como se ilustra en la estructura del proyecto, las interfaces del Nivel 2 se cargan bajo demanda de forma interna modificando dinámicamente la propiedad \texttt{PictureName} del objeto \texttt{PictureWindow} central configurado en el marco maestro estático \texttt{00\_Main\_Interface.pdl} \cite{1184, 1185}. La ventana inferior de avisos críticos de \textit{Alarm Logging} permanece fija en la zona del pie de página de forma persistente, garantizando que el operador no pierda la conciencia situacional ni la trazabilidad de eventos del nodo \texttt{WALM} de la planta durante el desplazamiento por el Nivel 2 \cite{1186, 1353, 1414}.